\def\ChapterCode{RSP}
\chapter
[\CO{[ATM]} Störungen der Atemwege und der Atmung]
{Störungen der Atemwege und der Atmung}
\label{topic:respiratorische-notfaelle}
\label{chp:atm}
\ChapterIntro
{%
~~~
}
{%
\begin{description}
  \item[Maintainer:] Sebastian Gabriel
  \item[Autoren:] Diverse
  \item[Reviewer:] Standard-Reviewprozess
  \item[Version:] {\DocVersion} ({\DocVersionInfo})
  \item[SHA1:] {\DocGitCommit}
\end{description}

{\TxTeilkapitelAASS}
}

% \begin{center}
%  \marginpar{\includegraphics[width=\linewidth]{./images/noauth/AASdev20090228-img75.png}
%  \AuthNotesPicLegalNotOk{img75}{Splash Pulmo Notfälle }}
% % 1-E
% \Layout{0}
% {\TxABCDE}
% {~
% \begin{ABCDEText}
%   \item[\ABCDE{1}] 
%   \item[\ABCDE{2}] 
%   \item[\ABCDE{3}] 
%   \item[\ABCDE{4}] 
%   \item[\ABCDE{A}] 
%   \item[\ABCDE{B}] 
%   \item[\ABCDE{C}] 
%   \item[\ABCDE{D}] 
%   \item[\ABCDE{E}] 
%   \item[\ABCDE{=}] 
%   \item[\ABCDE{\ldots}] 
% \end{ABCDEText}
% }
% {
% \begin{ABCDESlide}
%   \item[\ABCDE{Sz}] 
%   \item[\ABCDE{Ei}] 
%   \item[\ABCDE{Be}] 
%   \item[\ABCDE{Hb}] 
%   \item[\ABCDE{A}] 
%   \item[\ABCDE{B}] 
%   \item[\ABCDE{C}] 
%   \item[\ABCDE{D}] 
%   \item[\ABCDE{E}] 
%   \item[\ABCDE{=}] 
%   \item[\ABCDE{\ldots}] 
% \end{ABCDESlide}
% }
% {}
% {}
% {}
% 
% \Layout{0}
% {{\TxSAMPLER}}
% {
% \begin{SAMPLERText}
%   \item[\SAMPLER{S}] 
%   \item[\SAMPLER{A}] 
%   \item[\SAMPLER{M}] 
%   \item[\SAMPLER{P}] 
%   \item[\SAMPLER{L}] 
%   \item[\SAMPLER{E}] 
%   \item[\SAMPLER{R}] 
% \end{SAMPLERText}
% }
% {
% \begin{SAMPLERSlide}
%   \item[\SAMPLER{S}] 
%   \item[\SAMPLER{A}] 
%   \item[\SAMPLER{M}] 
%   \item[\SAMPLER{P}] 
%   \item[\SAMPLER{L}] 
%   \item[\SAMPLER{E}] 
%   \item[\SAMPLER{R}] 
% \end{SAMPLERSlide}
% }
% {}
% {}
% {}

\clearpage
\Layout{2}{{\TxQuerverweise}}{%
\begin{itemize}
  \item Lungenentzündung (Pneumonie): \dotfill \FullPrettyRef{topic:pneumonie}
  \item Tuberkulose:
  \dotfill \FullPrettyRef{topic:tuberkulose}
\end{itemize}
}{}{}{}{}





% 
% \section{Blutungen au dem Nasenraum}
% 
% % TODO: Nasenbluten
% 
% \Layout{2}
% {{\TxBeschreibung}}
% {}
% {}
% {}
% {}
% {}
% 


% \clearpage
\section{Mechanische Atemwegsverlegung}
\label{topic:mechanische-atemwegsverlegung}
\label{topic:fremdkoerperaspiration}

\Layout{60}{\TxBeschreibung}
{%
\index{Aspiration}%
Wenn ein Fremdkörper den Atemweg teilweise oder vollständig 
\E{blockiert} oder \E{behindert}, spricht man von einer
mechanische Verlegung. Abhängig von der Größe des
Fremdkörpers können sowohl die oberen als auch die unteren
Atemwege betroffen sein. Wenn der Fremdkörper in die
\E{unteren Atemwege} gelangt, spricht man von einer
(Fremdkörper-) \T{Aspiration}.

Nicht nur feste Körper, sondern auch Flüssigkeiten können
Auslöser einer Verlegung oder Aspiration sein. Besonders
Blut, Magensaft, Erbrochenes oder Speisebrei stellen
diesbezüglich eine Gefahr dar.
Besonders gefährdet hinsichtlich einer Atemwegsverlegung sind
Kinder, ältere Personen, bewusstseinseingschränkte Personen,
Patienten mit Schluckstörungen oder Allergiker (z.\,B.
Anschwellen der Schleimhäute nach einem Bienenstich).

Man unterscheidet zwischen einer \E{milden} (der Patient
kann noch sprechen und husten) und einer
\E{schweren Atemwegsverlegung} (Sprechen und Husten nicht
möglich) \citep{Erc2005Sec2,Erc2005Sec6,Erc2005DeSec2,Erc2005DeSec6}.



}
{%
\begin{itemize}
  \item Verlegung durch Fremdkörper
  \item Obere und untere Atemwege möglich
  \item Aspiration: untere Atemwege
  \item Besonders gefährdet:
  \begin{itemize}
    \item Kinder oder alte Personen
    \item Bewusstseinsgetrübte/-lose Personen
    \item Patienten mit Schluckstörungen
    \item Allergiker
  \end{itemize}
  \item Schweregrad:
  \begin{itemize}
    \item Milde Verlegung: \EE{Kann Sprechen \& Husten}
    \item Schwere Verlegung: Sprechen oder
    effektiv husten nicht möglich
  \end{itemize}
\end{itemize}
}
{./images/hirtler-aass/Fremdkoerper.pdf}
{Atemwegsverlegung (Schema). Der linke Bissen steckt in der Luftröhre fest, der
andere befindet sich in der Speiseröhre.}
{Hirtler.}




\Fazit{Die (Un-) Fähigkeit zu sprechen und zu husten ist das
wesentliche Kriterium um die Schwere einer Atemwegsverlegung
einzuschätzen.}

\Layout{0}
{\TxABCDE}
{~
\begin{ABCDEText}
  \item[\ABCDE{1}] Achte auf Essenreste, Spielsachen
  (Murmeln!) und andere potentielle Fremdkörper.
  %%%
  \item[\ABCDE{2}] Oft gestikuliert der Patient
  entsprechend. Bei fortgeschrittener Hypoxie kann der
  Patient {\RedFlag}\,zyanotisch sein.
  %%%
  \item[\ABCDE{3}] Infolge einer Hypoxie
  kann es zu einer {\RedFlag}\,\E{Bewusstseinstrübung} kommen.
  %%%
  \item[\ABCDE{4}] Das Leitsymptom ist Atemnot bzw.
  Fremdkörpergefühl, zusammen mit einer entsprechenden
  Anamnese.
  %%%
  \item[\ABCDE{A}] Der Atemweg ist teilweise oder ganz
  verlegt. Eventuell ist ein pfeifendes Atemgaeräusch
  (\I{Stridor}) zu hören. Bei einer hochgradigen
  oder kompletten Verlegung ist das Atemgeräusch
  abgeschwächt oder fehlt völlig. Der Fremdkörper ist oft
  nciht sichtbar, da er sich oft unterhalb des Rachenraums
  befindet. 
  %%%
  \item[\ABCDE{B}] Je nach Ausmaß der Verlegung kann die
  Sauerstoffsättigung erniedrigt sein. Der Patient weisst
  eine schnelle und erschwerte Atmung auf (\I{Tachypnoe}).
  \begin{itemize}
  \item \T[Atemwegsverlegung!milde]{Milde Verlegung}: Bei
  der milden Verlegung kann der Patient noch sprechen und effektiv husten. Er klagt
  meist über Atemnot und ein Fremdkörpergefühl, oft versucht er zu husten
  oder zu würgen. Oft ist auch ein pfeifendes Atemgeräusch
  während der Ein- und/oder Austamung zu hören
  (in-/exspiratorischer \E{Stridor}).
  \item \T[Atemwegsverlegung!schwere]{Schwere Verlegung}:
  Bei einer schweren Verlegung der Atemwege ist der Patient
  \EE{unfähig zu reden oder effektiv zu husten}. Eventuell kann
  er nur durch Nic\-k\-en oder Gesten kommunizieren. In besonders
  schweren Fällen kann es zu einem \EE{Atemstopp} und nach
  kurzer Zeit zu einer \EE{Bewusstlosigkeit}, gefolgt von einem
  \EE{Kreislaufstillstand} kommen.
\end{itemize}
  %%%
  \item[\ABCDE{C}] Der Patient ist tachykard.
  %%%
  \item[\ABCDE{D}] Oft liegt eine Hypertonie aufgrund des
  Stress vor. 
%   %%%
%   \item[\ABCDE{E}] 
%   %%%
%   \item[\ABCDE{\ldots}] 
%   %%%
  \item[\ABCDE{=}] Bei einer schweren Atemwegsverlegung ist
  der Patient vital bedroht. Bei der milden Verlegung ist
  die Einschätzung von den Befunden abhängig.
\end{ABCDEText}
}
{
\begin{ABCDESlide}
  \item[\ABCDE{1}]
  %%%
  \item[\ABCDE{2}]
  %%%
  \item[\ABCDE{3}]
  %%%
  \item[\ABCDE{4}] %
  \begin{itemize}
    \item Atemnot
    \item Fremdkörpergefühl
  \end{itemize}
  %%%
  \item[\ABCDE{A}] %
  \begin{itemize}
    \item Fremdkörper oft nicht sichtbar
    \item Evtl. pfeifendes Atemgeräusch (Stridor)
  \end{itemize}
  %%%
  \item[\ABCDE{B}] %
  \begin{itemize}
    \item Atemnot
    \item Schwere Verlegung:
    \begin{itemize}
      \item Unfähigkeit zu Sprechen oder effektiv zu husten
      \item Evtl. Bewusstlosigkeit → Kreislaufstillstand
    \end{itemize}
    \item Milde Verlegung:
    \begin{itemize}
      \item \EE{Kann Sprechen \& Husten}
    \end{itemize}
  \end{itemize}
  %%%
  \item[\ABCDE{C}]
  %%%
  \item[\ABCDE{D}]
  %%%
  \item[\ABCDE{E}]
  %%%
  \item[\ABCDE{\ldots}]
  %%%
  \item[\ABCDE{=}]
\end{ABCDESlide}
}
{}
{}
{}




\Layout{0}{Spezielle Techniken: Schläge zwischen die
Schulterblätter und Heimlich-Manöver}
{%
\begin{itemize}
  \item \T[Schulterblätter!Schläge zwischen die]{Schläge
  zwischen die Schulterblätter}:
  Durch kräftige Schläge auf den Rücken zwischen die
  Schulterblätter kann  ein festsitzender Fremdkörper evtl.
  gelockert und ausgehustet werden.
  \item \T{Heimlich-Manöver}:
  \label{topic:heimlich-manoever}%
  \Lang{Syn.} \I{Heimlich-Handgriff}. Beim Heimlich-Manöver
  versucht man durch Kompression des Bauchraums den Fremdkörper
  durch den entstandenen Überdruck aus den Atemwegen zu
  entfernen.
  
  Der Helfer umfasst von hinten den Oberbauch des Patienten und
  bildet mit der einen Hand eine Faust und umschliesst diese
  mit der anderen Hand. Nun legt er die Faust unterhalb der
  Rippen und des Brustbeins an und zieht sie dann ruckartig
  gerade nach hinten zu sich.
  
  \Ca{Verletzungsgefahr!} Durch Anwendung des
  Heimlich-Handgriffes kann es zu schweren inneren Verletzungen
  kommen (Magen, Leber, \ldots), die Anwendung ist daher auf
  bedrohliche Situationen beschränkt.
\end{itemize}
}{%
\begin{itemize}
  \item Schläge zwischen Schulterblätter
  \item Heimlich-Manöver
  \begin{itemize}
    \item Kompression des Bauchraumes → Überdruck
    \item Ansatzpunkt unterhalb des Rippenbogens, ruckartige Bewegung
    \item \Ca{Verletzungsgefahr!}
  \end{itemize}
\end{itemize}
}
{}
{}
{}


\begin{PictureSeries}{}{Bilderserie: Atemwegsverlegung}
\PictureSeriesRowWide{2}
{./images/hirtler-aass/Fremdkoerper.pdf}
{Atemwegsverlegung (Schema). Der linke Bissen steckt in der Luftröhre fest, der
andere befindet sich in der Speiseröhre. \SourcePicture{Hirtler.}}
{./images/wikipedia-pd/Abdominal_thrusts3.jpg}
{Heimlich-Manöver \SourcePicture{U.\,S. Naval Hospital, Okinawa, Japanese intern Naoko Uehara teaches
proper technique for treating a choking victim with the assistance of Hospital
Corpsman Second Class Ross Francis. Amanda M. Woodhead, USMC / WmCo, PD}}
{empty}
{}
{empty}
{}
\end{PictureSeries}


\MassnahmenMK{Spezielle Maßnahmen}{Mechanische
Atemwegsverlegung}{m:atemwegsverlegung} {%
\EE{Taktik}: 
\\\Speech{Schwer: \Ca{Vitale Bedrohung!} Zeitkritisch; wenn
möglich rasche Entfernung des Fremdkörpers. 
\\
Mild: Symptomatische
Therapie}
}{}

 

% \Layout{2}{\dsliterary}{%
\nocite{Wuttig2009}
% }{}{}{}{}


\section{Asthma bronchiale}
\label{topic:asthma-bronchiale}

\Layout{60}{\TxBeschreibung{}}
{
\DefinitionInPara{Das \T{Asthma bronchiale} ist eine chronische 
Erkrankung mit Episoden von Anfällen mit massiver Atemnot oder
Atembehinderung durch Verengung der Bronchien.}
Im akuten Asthmaanfall kommt es zu einer Verengung der
Bronchien und vermehrter Schleimbildung, die \EE{Ausatmung}
wird massiv erschwert und der Patient hat das Gefühl nicht
genug Luft zu bekommen.
}{%
Wiederholte Anfälle von Atemnot durch
\begin{itemize}
  \item Verengte Bronchien
  \item Vermehrte Schleimbildung
  \item Behinderte Ausatmung
\end{itemize}
}
{./images/asthma-paramedic.jpg}
{Patientin mit einem akuten Asthmaanfall. Sie sitzt auf einer Treppe und stützt sich nach hinten mit
den Armen ab. Die Erstmaßnahmen bei vital bedrohten
Patienten wurden ergriffen: Situationsrechte Lagerung,
Sauerstoffgabe, Reanimationsbereitschaft, Notarztnachforderung und Monitoring.\label{fig:asthma-paramedic}
}
{Gabriel}

Asthma ist bei den meisten Patienten oft schon bekannt und
kann bzw. muss erfragt werden. Viele Patienten haben auch
einen oder mehrere Sprays oder Inhalationsgeräte zur Dauer-
oder Akutbehandlung. 
Besonders betroffen vom Asthma bronchiale sind Patienten
mit bekannten Allergien. Folglich kann alles, was bei dem
Patienten eine Allergie verursacht auch einen Asthmaanfall
auslösen (Pollen, Gräser, Katzen, Zigarettenrauch, \ldots). 

Ist bei Eintreffen des Rettungsdienstes noch keine
Besserung der Atemnot eingetreten, ist diese Atemnot als
massiv und lebensbedrohend einzuschätzen. Es ist nicht
damit zu rechnen, dass der Zustand von alleine wieder
vergeht und bedarf einer medikamentösen Behandlung durch einen
Notarzt. 






\Layout{0}
{\TxABCDE}
{~
\begin{ABCDEText}
  \item[\ABCDE{1}] Kalte Temperaturen oder Allergene
  (Gräser, Pollen, \ldots) begünstigen einen Asthma-Anfall.
  %%%
  \item[\ABCDE{2}] Der Patient nimmt meist automatisch und unbewusst eine
  Position ein, die seine Atemhilfsmuskulatur unterstützt: Er
  sitzt aufrecht und stützt sich oft nach vorne oder hinten
  ab.
  %%%
  \item[\ABCDE{3}] Im schweren Anfall kann es zu
  {\RedFlag}\,\E{Bewusstseinsstörungen} kommen.
  %%%
  \item[\ABCDE{4}] Das Leitsymptom ist die Atemnot.
  %   %%%
  %   \item[\ABCDE{A}]
  %%%
  \item[\ABCDE{B}] Die Atemfrequenz ist erhöht
  (\E{Tachypnoe}), die {\RedFlag}\,\E{Sauerstoffsättigung}
  kann erniedrigt sein. Bei einer manifesten Hypoxie kann eine
  {\RedFlag}\,\E{Zyanose} beobachtbar sein.
  
  Oft ist ein pfeifendes,
  brummendes oder giemendes Atemgeräusch, besonders während
  der Ausatmung (\E{Exspiration}, \E{exspiratorischer
  Stridor}) wahrnehmbar. Die Exspiration ist ausserdem
  verlängert. Trockener Husten ist häufig.
  
  \Z{Unterschätze nie die Schwere eines Asthmaanfalles!
  Gefährlich sind besonders die \protect\enquote*{ruhigen},
  wenig klagenden Patienten, diese haben oft ein anderes
  Luftnot-Empfinden. Diese Patienten sind Hochrisiko-Patienten
  und haben auch mit der \ce{O2}-Gabe Probleme! \ce{O2}-Gabe nur
  bis max. 95\PC* Sauerstofsättigung!
  \citep{Olschewski2009}\A{\footnote{Vortrag Univ-Prof.Dr. Horst Olschewski, 5. AGN-Kongress 2009; Notfall Rettungsmed (2009) 12:322}}}
  
  Im \EE{schweren Anfall} ist oft \E{kein}, oder nur mehr ein
  \E{leises Atemgeräusch} hörbar, da kaum mehr Luft bewegt
  wird \Z{(\I{silent chest})}. Die
  {\RedFlag}\,Sauerstoffsättigung ist stark vermindert und der
  Patient ist {\RedFlag}\,\E{zyanotisch}. Im Endstadium kann
  die {\RedFlag}\,\E{Atemarbeit erschöpft} sein.
  %%%
  \item[\ABCDE{C}] Die Herzfrequenz ist meist erhöht
  (\E{Tachykardie}).
  
  
  %%%
  \item[\ABCDE{D}]
  %%%
  \item[\ABCDE{E}]
  %%%
  \item[\ABCDE{\ldots}]
  %%%
  \item[\ABCDE{=}]
\end{ABCDEText}
}
{
\begin{ABCDESlide}
  \item[\ABCDE{1}] 
  %%%
  \item[\ABCDE{2}] Einsatz der Atemhilfsmuskulatur
  %%%
  \item[\ABCDE{3}] Schwerer Anfall:
  {\RedFlag}\,Bewusstseinsstörungen
  %%%
  \item[\ABCDE{4}] Akute Atemnot
%   %%%
%   \item[\ABCDE{A}] 
  %%%
  \item[\ABCDE{B}] Akute Atemnot mit trockenem
  Husten, Verlängerte Ausatmungsphase, AF $\uparrow$,
  Exspiratorisches Atemgeräusch (Stridor)
  
  Schwerer Anfall: Kein oder leises Atemgeräusch,
  {\RedFlag}\,Sp\ce{O2} $\downarrow\downarrow$
 
  Endstadium: {\RedFlag}\,Atemarbeit erschöpft
  %%%
  \item[\ABCDE{C}] HF $\uparrow$
  
  Schwerer Anfall: HF $\downarrow$
  %%%
  \item[\ABCDE{D}] 
  Schwerer Anfall: RR $\downarrow$
  %%%
  \item[\ABCDE{E}] 
  %%%
  \item[\ABCDE{\ldots}] 
  %%%
  \item[\ABCDE{=}] 
\end{ABCDESlide}
}
{}
{}
{}

\subparagraph{\Ca{Schwerer Anfall}}




\begin{itemize}
  \item 
  \item
  \item 
  \item 
  \item 
  \item Schwerer Anfall:
  \begin{itemize}
    \item 
    \item 
    \item 
    \item 
  \end{itemize}
\end{itemize}


\Layout{0}
{\TxABCDE}
{~
\begin{ABCDEText}
  \item[\ABCDE{1}]
  \item[\ABCDE{2}] 
  \item[\ABCDE{3}]
  \item[\ABCDE{4}] 
  \item[\ABCDE{A}] 
  \item[\ABCDE{B}] 
  \item[\ABCDE{C}] 
  \item[\ABCDE{D}] 
  \item[\ABCDE{E}] 
  \item[\ABCDE{=}] 
\end{ABCDEText}
}
{
\begin{ABCDESlide}
  \item[\ABCDE{Sz}] 
  \item[\ABCDE{Ei}] 
  \item[\ABCDE{Be}] 
  \item[\ABCDE{Hb}] 
  \item[\ABCDE{A}] 
  \item[\ABCDE{B}] 
  \item[\ABCDE{C}] 
  \item[\ABCDE{D}] 
  \item[\ABCDE{E}] 
  \item[\ABCDE{=}] 
\end{ABCDESlide}
}
{}
{}
{}

\Layout{0}
{{\TxSAMPLER}}
{
\begin{SAMPLERText}
  \item[\SAMPLER{S}] 
  \item[\SAMPLER{A}]
  \item[\SAMPLER{M}] 
  \item[\SAMPLER{P}]
  \item[\SAMPLER{L}]
  \item[\SAMPLER{E}] 
  \item[\SAMPLER{R}]
\end{SAMPLERText}
}
{
\begin{SAMPLERSlide}
  \item[\SAMPLER{S}]
  \item[\SAMPLER{A}]
  \item[\SAMPLER{M}] 
  \item[\SAMPLER{P}] 
  \item[\SAMPLER{L}]
  \item[\SAMPLER{E}] 
  \item[\SAMPLER{R}]
\end{SAMPLERSlide}
}
{}
{}
{}

\Layout{0}{Status Asthmaticus}{%
\index{Status!asthmaticus}
\label{topic:status-asthmaticus}
Der
\protect\enquote{Status} ist eine gefürchtete Komplikation des
Asthmaanfalles bei der die Atemnot nicht wieder vergeht. Es
besteht die Gefahr, dass der Patient nach einiger Zeit keine
Kraft mehr zum Atmen hat (Erschöpfung). Der Status
asthmaticus kann tödlich enden.
}{
\begin{itemize}
  \item Lebensgefährliche Komplikation:
  \item Anfall hört nicht mehr auf
\end{itemize}
}{}{}{}




% \begin{figure}[h]
%     \includegraphics{./images/copd-bronchus.png}
%     
%     \Captionof{figure}{Bronchus bei COPD.
%     \AuthNotesPicLegalNotOk{copd-bronchus}{}}
% \end{figure}

\MassnahmenMK{Spezielle Maßnahmen}{Akuter Asthmaanfall}
{m:asthmaanfall}{
\EE{Taktik}: \Speech{Linderung der Atemnot und rasche
medikamentöse Therapie \Ca{Vitale Bedrohung im Falle
bleibender schwerer Atemnot bzw. Status asthmaticus}
}
 }{}

% \MassnahmenNG{Spezielle Maßnahmen}{Akuter Asthmaanfall}
% {\label{m:asthmaanfall}}{
% \EE{Taktik}: \Speech{Linderung der Atemnot und rasche
% medikamentöse Therapie \Ca{Vitale Bedrohung im Falle
% bleibender schwerer Atemnot bzw. Status asthmaticus}
% }
% 
% \begin{itemize}
%   \item Lagerung: \EE{Oberkörper hoch}, abstützen lassen
%   (Unterstützung der Atemhilfsmuskulatur)
%   \item \ce{O2}-Gabe gemäß
%   \prettyref{m:sauerstoffberieselung} 
%   
%   {Achtung: In seltenen Fällen
%   kann es vorkommen, dass bei hochdosierter
%   Sauerstoffgabe der Patient weniger atmet und
%   Bewusstseinsstörungen auftreten (\ce{CO2}-Narkose)! Siehe
%   \protect\enquote{\Abk{COPD}} (\prettyref{topic:copd}). Bei sorgfältiger
%   Überwachung des Patienten stellt die \ce{O2}-Gabe in der
%   Praxis jedoch kein Problem
%   dar \citep{ScholzNotfallmedizin2}.}
%   \hfill\Commentary[Asthma, Sauerstoffgabe]{Asthma:
%   \protect\ce{O2} bis eine Sauerstoffsättigung von \VonBis{94}{98}\PC* erreicht ist, oder 8\LPerMin*, gem.
%   \cite{AsboeLehrmeinungRd2011-11}. Die einschlägige
%   Literatur ist widersprüchlich. Besonderes betont muss die Notwendigkeit zur Überwachung des
%   Patienten bezüglich Bewusstseinsstörungen und Atmung
%   werden!)}
%   \item Beruhigen, voratmen, durch fast
%   geschlossene Lippen ausatmen lassen (Lippenbremse)
%   \item Sprays, die Patient evtl. bei sich hat, dürfen
%   nicht mehr genommen werden.
%   \hfill\Commentary[Asthma / Spray]{Oft stellt sich die
%   Frage \protect\enquote{Darf der Patient seinen (vom Arzt vielleicht gerade für
%   diese Situation verschriebenen) Spray nehmen?}
%   Grundsätzlich soll beim schweren Asthma-Anfall davon
%   abgesehen werden, \protect\enquote{blind} ein inhalatives Beta-Mimetikum
%   zu geben, das zweifelhaft ist, ob die Luftpassage
%   ausreicht, umd den Wirkstoff in die Lunge zu befördern.
%   Eine generelle Freigabe, verschriebene Srpays einzunehmen,
%   kann in Anbetracht der Nebenwirkungen bei
%   Schleimhautresorption nicht gegeben werden.}
% \end{itemize}
% 
% \cite{AsboeLehrmeinungRd2011-11}
%  }

\bigskip







% \clearpage % TEMP temp clearpage
\section{Chronische Bronchitis und 
Chronische Obstruktive Lungenerkrankung
(COPD)}
\label{topic:copd}
\nocite{Renz-Polster2006} 

% Renz-Polster2006 p442

% \HeadingSub{}

\DefinitionInPara{Die \T{chronische Bronchitis} ist
eine chronische entzündliche Schleimhautschädigung  der unteren Atemwege.}
\DefinitionInPara{Die \T{COPD}
(\T[Lungenerkraknkung!chronisch-obstruktive]{Chronische Obstruktive
Lungenerkrankung}\ZFootnote{\Lang{engl.} Chronic
obstructive pulmonary disease\index{Chronic
obstructive pulmonary disease|Z})} ist eine chronische entzündliche
Schleimhautschädigung, welche eine zunehmende obstruktive 
Atemwegseinschränkung aufweist.}
Am Anfang steht die chronische Bronchitis, welche durch
Husten mit schleimigem Auswurf gekennzeichnet ist
({\I{Raucherhusten}}). Es kommt dabei zu einer gesteigerten
Entzündungsantwort auf eingeatmete Stoffe (Zigarettenrauch,
Umweltschadstoffe, \ldots). Wenn man in der
Lungenfunktionsuntersuchung eine Atemwegseinschränkung
nachweisen kann, spricht man von der chronisch-obstruktiven
Lungenerkrankung (\Abk{COPD}). 
Diese geht mit bleibenden Veränderungen der unteren Atemwege einher, 
diese sind in der Grafik  \prettyref{fig:copd-veraenderungen} dargestellt.


\begin{figure}[H]

\includegraphics[width=\linewidth]{./images/hirtler-aass/Alveolen-Bronchus-COPD-edited2.png}
\Captionof{figure}{\label{fig:copd-veraenderungen}Veränderung der
Atemwege bei der COPD. \TabSub{Links:} Schema eines gesunden
Bronchus und einer gesunden Alveole. \TabSub{Oben rechts:} Bei der COPD sind
die kleinen Luftwege verschleimt und verengt. \TabSub{Unten rechts:} Die
Lungenbläschen (Alveolen) sind überbläht, weil die Luft nur erschwert wieder entweichen kann.
\SourcePicture{Hirtler.}}
\end{figure}



\Layout{0}{\TxSymptome}{%
Die \Abk{COPD} ist durch eine voranschreitende
Verschlechterung der Atemleistung gekennzeichnet, erstes
Anzeichen ist ein Husten mit Auswurf, der klassische
\protect\enquote{Raucherhusten}.
Je nach Schweregrad kommt es zu Zeichen der
Atemwegsverlegung (Obstruktion) und Ateminsuffizienz: Die
\E{Ausatmung} (Exspiration) ist erschwert, der Patient klagt
über \E{Belastungsdyspnoe}, ein Engegefühl, \E{Husten} und
hat Tachypnoe und evtl. Zyanose. Man kann oft ein
\E{brummendes, pfeifendes oder giemendes Atemgeräusch}
hören.

\Z{Durch die erschwerte Ausatmung kommt es zu einer
chronischen Überblähung der Lungenbläschen und zu einer \protect\enquote{Faß-förmigen}
Verformung des Brustkorbes. Im Endstadium zeigen sich
Zeichen einer Rechtsherzinsuffizienz
(\prettyref{topic:rechtsherzinsuffizienz}) aufgrund einer
Störung im Lungenkreislauf.}

Die \EE{Sauerstoffsättigung} bei einem COPD-Patienten ist
auch ohne akute Symptome meist \E{deutlich vermindert}
($\sim$\,{90}\PC*). Der sonst übliche Normalwert
gilt hier nicht! Bei fortgeschrittener Ateminsuffizienz
bekommt der Patient oft eine \E{Heimsauerstofftherapie} verordnet.
}{
\begin{itemize}
  \item Husten mit Auswurf, anfänglich
  \protect\enquote*{Raucherhusten}
  \item Voranschreitende, sich verschlechternde \E{Ateminsuffizienz}:
  \begin{compactitem}
    \item Belastungsdyspnoe (bei schwererer \Abk{COPD}
    auch in Ruhe)
    \item Ausatmung (Exspiration) erschwert
    \item Brummendes, pfeifendes od. giemendes Ausatemgeräusch
    \item Evtl. chronische Zyanose
    \item Evtl. Heimsauerstoff
    \item Sp\ce{O2} chronisch niedrig
  \end{compactitem}
\end{itemize}
}{}{}{}



\Layout{0}{Exazerbationen}{%
Der \Abk{COPD}-Patient erleidet öfters plötzliche
Verschlechterungen, sog. \E{Exazerbationen}. Diese sind
meist durch Infekte ausgelöst und durch vermehrte Atemnot,
Husten und Auswurf gekennzeichnet.
}{
\begin{itemize}
  \item Plötzliche Verschlechterungen
  \item Meist Infekt-bedingt
  \item Vermehrte Atemnot,
  Husten und Auswurf
\end{itemize}
}{}{}{}


% 
% \subparagraph*{Sammelbegriff für}

% \begin{itemize}
%     \item Lungenemphysem (Überblähung der Lungenbläschen)
%     \item chron. obstruktive Bronchitis
% \end{itemize}








\begin{figure}[h]
\begin{center}
% \includegraphics{./images/copd-stufen.png}
\includegraphics[width=\linewidth]{./images/hirtler-aass/COPD-Stufen_edited.pdf}

\Captionof{figure}{Eine COPD
entsteht nicht plötzlich: Ein COPD{\textquotesingle}ler hat eine
\protect\enquote{Karriere} hinter sich. \SourcePicture{Hirtler}}
\end{center}
\end{figure}


\Layout{0}{Probleme mit Sauerstoff bei
COPD-Patienten}
{
\index{Sauerstoff!COPD-Patient}
Normalerweise wird der Atemantrieb über die
\ce{CO2}-Konzentration im Blut gesteuert. Die
\ce{CO2}-Konzentration ist bei \Abk{COPD}-Patienten jedoch chronisch
hoch, der Körper gewöhnt sich an diesen Zustand. Er
regelt dann den Atemantrieb direkt über den \ce{O2}-Gehalt im Blut.
Wird der \ce{O2}-Gehalt plötzlich künstlich erhöht, fehlt
der \index{COPD!Atemantrieb}Atemantrieb und der Patient kann
\EE{Bewusstseinsstörungen} (\protect\enquote{\ce{CO2}-Narkose}) und einen
\EE{Atemstillstand} erleiden.

}{
\Ca{Achtung!}

\SymbolText
}{}{}{}

\MassnahmenMK{Spezielle Maßnahmen}{COPD-Exazerbation}
{m:copd-exzerbation}
{
\EE{Taktik}: \Speech{Linderung der Atemnot und rasche
medikamentöse Therapie \Ca{Vitale Bedrohung bei anhaltender
schwerer Atemnot!}}
}{}


% \MassnahmenNG{Spezielle Maßnahmen}{COPD-Exazerbation}{}{
% \EE{Taktik}: \Speech{Linderung der Atemnot und rasche
% medikamentöse Therapie \Ca{Vitale Bedrohung bei anhaltender
% schwerer Atemnot!}}
% 
% \begin{itemize}
%   \item
%   \TxBeiVit \TxMassVit 
%   \item \EE{\ce{O2}-Gabe:} \Ca{Ausnahme!}
%   
%   \E{Sauerstoff vorsichtig dosieren,
%   anfänglich nur \EE{\VonBis{2}{3}{\LPerMin*}}. Wenn der
%   Patient bereits Heimsauerstoff benutzt, 
%   \VonBis{1}{2}\LPerMin* höher zu dosieren.   
%   Heimsauerstoff soll jedenfalls weiter gegeben werden! }
%   
%   \EE{Atmung
%   und Bewusstsein müssen bei der \ce{O2}-Gabe besonders
%   sorgfältig überwacht werden!}  
%   \item Lagerung: \EE{Oberkörper hoch}
%   \item Voratmen, \I{Lippenbremse}
%   \item Eigener Spray darf nicht mehr genommen werden
%   
% \end{itemize}
% \AuthNotes{Spray bei COPD -- Regelung wie bei Nitro??}
% }

% \Layout{2}{\Z{Literatur}}
% {
% \Z{
% \begin{inparadesc}
%   \item[Kasuistiken:] \citep{Knacke200601}
% \end{inparadesc}
% }
% }{}{}{}{}

\nocite{Knacke200601}


\section{Lungenembolie}
\label{topic:lungenembolie}



\Layout{60}{{\TxBeschreibung} und Ursachen}
{%
\DefinitionInPara{Als \T{Lungenembolie} (\I{Pulmonalembolie}) wird
der \EE{Verschluss einer Lungenarterie}, oft  durch Einschwemmen eines 
\EE{Gerinnsels} (\I{Embolus}), bezeichnet.}
Der Embolus, der die Lungenarterien verstopft, entsteht
normalerweise nicht in den Lungengefäßen selbst, sondern er
wird über die Blutbahn angeschwemmt. Sehr oft löst sich ein Thrombus bei
einer \EE{tiefen Beinvenenthrombose} und schwimmt im Blut bis in
die Lunge. Daher gelten die
Risikofaktoren der tiefen Beinvenenthrombose
(Immobilisation, lange Reisen, \ldots) auch bei der
Lungenembolie. 
Auch bei \EE{Vorhofflimmern} können sich Gerinnsel im 
rechten Vorhof des Herzens bilden und eine Embolie verursachen.
Ebenso können
Gerinnungstörungen oder Krebserkrankungen zur Bildung von 
Blutgerinnseln führen.
}{
\begin{itemize}
  \item Verschluss einer Lungenarterie
  \item Oft eingeschwemmt, vor allem:
  \begin{itemize}
    \item Tiefe Beinvenenthrombose
    \item Rechter Vorhof bei Vorhofflimmern
  \end{itemize}
\end{itemize}
}
{./images/hirtler-aass/Thrombus-Embolie.pdf}
{Herkunft des Thrombus.}
{Hirtler.}




\Layout{0}
{\TxABCDE}
{~
\begin{ABCDEText}
  \item[\ABCDE{Szeneüberblick}]
  \item[\ABCDE{Eindruck}] 
  \item[\ABCDE{Bewusstsein}]
  \item[\ABCDE{Hauptbeschwerde}] Das Leitsymptom ist die \EE{Atemnot}, eventuell in Kombination mit (atemabhängigen) Thorax- oder Rückenschmerzen. 
  \item[\ABCDE{A}] 
  \item[\ABCDE{B}] Meist atmet der Patient schnell (Tachypnoe) und klagt über
einen \EE{atemabhängigen Thoraxschmerz}.
  \item[\ABCDE{C}]  Daneben kommt es zu
einer \E{Tachykardie} und Angstgefühlen., oft findet man deutlich hervortretende, \EE{gestaute Venen}
am Hals aufgrund des Blutrückstaus. Evtl. kann man auch einen
\E{arrhythmischen Puls} aufgrund eines Vorhofflimmerns, ertasten.
  \item[\ABCDE{D}] 
  \item[\ABCDE{E}] 
  \item[\ABCDE{\ldots}] Je nach Ursache der Embolie und Herkunft des Embolus kann
man Zeichen anderer Krankheitsbilder wahrnehmen: Bei einer
\EE{tiefen Beinvenenthrombose} ist ein Bein geschwollen,
überwärmt, rot/rot-bläulich verfärbt und schmerzhaft
(\prettyref{topic:thrombose}). 

Eine Lungenembolie kann man leicht mit einem Herzinfarkt
oder Angina pectoris verwechseln. Dies hat aber bei der
Erstversorgung kaum Konsequenz, die Erstversorgen ist gleich.
  \item[\ABCDE{=}] 
\end{ABCDEText}
}
{
\begin{ABCDESlide}
  \item[\ABCDE{Sz}] 
  \item[\ABCDE{Ei}] 
  \item[\ABCDE{Be}] 
  \item[\ABCDE{Hb}] Dyspnoe, (atemabhängige) Thorax- od. Rückenschmerzen, Angstgefühl
  \item[\ABCDE{A}] 
  \item[\ABCDE{B}] Atemnot, Tachypnoe
  \item[\ABCDE{C}] Tachkardie, evtl. gestaute Halsvenen (Blutrückstau) , evtl. arrhythmischer Puls bei Vorhofflimmern

  \item[\ABCDE{D}] 
  \item[\ABCDE{E}] 
  Evtl. Bein: einseitig geschwollen,
  überwärmt, rot/rot-bläulich verfärbt, schmerzhaft (Beinvenenthombose)
  \item Evtl. arrhythmischer Puls (Vorhofflimmern)
  \item[\ABCDE{=}] 
\end{ABCDESlide}
}
{}
{}
{}

\Layout{0}
{{\TxSAMPLER}}
{
\begin{SAMPLERText}
  \item[\SAMPLER{S}] 
  \item[\SAMPLER{A}]
  \item[\SAMPLER{M}] 
  \item[\SAMPLER{P}]
Oft ist bei den Patienten auch \EE{Vorhofflimmern} als
chronische Erkrankung bekannt.zu leiden.
  \item[\SAMPLER{L}]
  \item[\SAMPLER{E}] 
  \item[\SAMPLER{R}] \EE{Immobilisation}, z.\,B. durch Bettlägrigkeit,
Liegegipsverbände oder lange Reisen, innerhalb der letzten 4
Wochen, erhöht das Risiko einer Thrombose und damit auch das
einer Lungenembolie. \Z{Weiters tritt die Lungenembolie gehäuft in der Schwangerschaft
auf. Rauchen sowie die Einnahme der \protect\enquote{Pille} erhöhen
ebenso das Risiko.}
\end{SAMPLERText}
}
{
\begin{SAMPLERSlide}
  \item[\SAMPLER{S}]
  \item[\SAMPLER{A}]
  \item[\SAMPLER{M}] 
  \item[\SAMPLER{P}] 
  \item[\SAMPLER{L}]
  \item[\SAMPLER{E}] 
  \item[\SAMPLER{R}]Immobilisation, Vorhofflimmern, \Z{Schwangerschaft}, \Z{Rauchen}, \Z{Verhütungspille}
\end{SAMPLERSlide}
}
{}
{}
{}





\MassnahmenMK{Spezielle
Maßnahmen}{Lungenembolie}{m:lungenembolie}{ \EE{Taktik}: \Speech{\Ca{Vitale Bedrohung!} Atemnot lindern,
rasche medikamentöse Therapie}
}{}

% \MassnahmenNG{Spezielle Maßnahmen}{Lungenembolie}{}{
% \EE{Taktik}: \Speech{\Ca{Vitale Bedrohung!} Atemnot lindern,
% rasche medikamentöse Therapie, ärztliche Therapie sinnvoll}
% 
% \begin{itemize}
%   \item {\TxMassVitBes}
%   \begin{itemize}
%     \item Lagerung: Oberkörper hoch
%     \item \ce{O2}-Gabe großzügig
%   \end{itemize}
%   \item Striktes Bewegungsverbot!
%   \item Beengende Kleidung öffnen
% \end{itemize}
% }


\section{Lungenödem}
\label{topic:lungenoedem}

\nopagebreak

\Layout{0}{{\TxBeschreibung} und Ursachen}
{%
\DefinitionInPara{Als \T{Lungenödem} wird die 
Flüssigkeitsansammlung im Zwischengewebsraum und in den
Alveolen und der Lungen bezeichnet.}
Ein Lungenödem kann entweder aufgrund einer
\EE{Herzschwäche} (\E{kardial}) oder aufgrund
  von \EE{Inhalation giftiger Substanzen} (\E{toxisch})
entstehen:
\begin{itemize}
  \item Beim \EE{kardialen Lungenödem} versagt die Auswurfleistung
  des linken Herzens (\EE{Linksherzinsuffizienz}), dadurch
  kommt es zu einem \EE{Rückstau} des Blutes in die Lunge.
  Aufgrund der Druckerhöhung tritt Flüssigkeit in das Gewebe
  über, es bilden sich Ödeme im Gewebe bzw.
  \EE{Flüssigkeitsansammlungen} in den Alveolen. Auslöser kann
  zum Beispiel ein Herzinfarkt sein, der die Funktion des
  Hermuskels massiv in Mitleidenschaft zieht.
  \item Beim \EE{toxischen  Lungenödem} kommt es aufgrund reizender
  Stoffe zu einer \EE{Entzündung}, und damit zu einem
  \EE{Flüssigkeitsaustritt} im Lungengewebe.
\end{itemize}

In beiden Fällen behindert die Flüssigkeit den
\EE{Gasaustausch} massiv, beim kardialen Lungenödem kommt
noch das begleitende Herzproblem dazu.

}{
\begin{itemize}
  \item Flüssigkeitsansammlung im Zwischengewebsraum und in 
  Alveolen der Lunge
  \item Arten
  \begin{itemize}
    \item \EE{Kardial}: (Links-)Herzinsuffizienz (Blutrückstau); 
    z.\,B. bei MCI, zusätzlicher Belastung bei vorgeschädigtem Herz
    (Infekt), \ldots
    \item \EE{Toxisch}: Inhalation giftiger Substanzen
  \end{itemize}
  \item Gasaustausch beeinträchtigt
\end{itemize}

}{}{}{}









\Layout{0}{\TxSymptome}
{%
\index{Brodelnde Atemgeräusche}%
Das Leitsymptom ist das \EE{brodelnde Atemgeräusch}, welches
in schweren Fällen bereits ohne Hilfsmittel hörbar und
typisch für das Lungenödem ist. Es klingt, als würde jemand
mit einem Strohhalm in ein Glas Wasser hineinblasen. Im
Extremfall kann man es bereits \E{auf mehrere Meter
Entfernung} wahrnehmen.
Weitere häufige Symptome sind \EE{Atemnot} und \EE{Zyanose}.
Der Patient bekommt oft nur \E{aufrecht sitzend} genügend
Luft (Orthopnoe).
}{
\begin{compactitem}
  \item Brodelndes Atemgeräusch
  \item Atemnot
  \item Zyanose
  \item Patient bekommt nur im Sitzen Luft
\end{compactitem}
}{}{}{}

\Layout{0}
{\TxABCDE}
{~
\begin{ABCDEText}
  \item[\ABCDE{Szeneüberblick}]
  \item[\ABCDE{Eindruck}]
  \item[\ABCDE{Bewusstsein}]
  \item[\ABCDE{Hauptbeschwerde}]
  \item[\ABCDE{A}]
  \item[\ABCDE{B}]
  \item[\ABCDE{C}]
  \item[\ABCDE{D}]
  \item[\ABCDE{E}]
  \item[\ABCDE{=}] Patienten mit einem deutlich hörbaren Lungenödem (v.\,a. bei
  kardialer Ursache) sind schwer und \EE{lebensbedrohlich
  krank}. Das Herz verbraucht oft schon die allerletzten Reserven. In
  schweren Fällen kann schon die Aufregung durch den Transport
  im Stiegenhaus zum Auto zu einem Herz-Kreislaufstillstand
  führen.
  
  Es kommt tatsächlich relativ häufig vor, dass diese
  Patienten \EE{im Stiegenhaus} oder \EE{im Fahrzeug}
  \EE{reanimationspflichtig} werden. \Ca{Solche Patienten müssen
  daher \EE{an Ort und Stelle} eine entsprechende (ärztliche)
  Therapie erhalten bevor sie transportiert werden können.}
  
  Bei leichten Formen des Lungenödems sind die
  Atemgeräusche nicht mit freiem Ohr, sondern
  nur mittels Stethoskop zu hören.  Die Patienten zeigen auch
  sonst keine Alarmsymptome, sondern nur Ödeme und eine mäßige
  Kurzatmigkeit. Diese Patienten sind i.\,d.\,R. nicht
  unmittelbar vital bedroht, ein bestmögliches Monitoring ist jedoch wichtig.
\end{ABCDEText}
}
{
\begin{ABCDESlide}
  \item[\ABCDE{Sz}]
  \item[\ABCDE{Ei}]
  \item[\ABCDE{Be}]
  \item[\ABCDE{Hb}]
  \item[\ABCDE{A}]
  \item[\ABCDE{B}]
  \item[\ABCDE{C}]
  \item[\ABCDE{D}]
  \item[\ABCDE{E}]
  \item[\ABCDE{=}]
  \begin{itemize}
    \item Hörbares Lungenödem: \Ca{Lebensgefahr!}
    \item Leichte Form: Keine Alarmsymptome
  \end{itemize}
\end{ABCDESlide}
}
{}
{}
{}

\Layout{0}
{{\TxSAMPLER}}
{
\begin{SAMPLERText}
  \item[\SAMPLER{S}] 
  \item[\SAMPLER{A}]
  \item[\SAMPLER{M}] 
  \item[\SAMPLER{P}]
  \item[\SAMPLER{L}]
  \item[\SAMPLER{E}] 
  \item[\SAMPLER{R}]
\end{SAMPLERText}
}
{
\begin{SAMPLERSlide}
  \item[\SAMPLER{S}]
  \item[\SAMPLER{A}]
  \item[\SAMPLER{M}] 
  \item[\SAMPLER{P}] 
  \item[\SAMPLER{L}]
  \item[\SAMPLER{E}] 
  \item[\SAMPLER{R}]
\end{SAMPLERSlide}
}
{}
{}
{}

\Layout{0}{Einschätzung}{%


}{
}{}{}{}

\MassnahmenMK{Spezielle Maßnahmen}{Lungenödem,
leicht}{m:lungenoedem-leicht}
{ \EE{Taktik}:
\Speech{Belastung minimieren, bestmögliches Monitoring}
}{}
% \MassnahmenNG{Spezielle Maßnahmen}{Lungenödem,
% leicht}{\label{m:lungenoedem-leicht}}
% { \EE{Taktik}:
% \Speech{Belastung minimieren, bestmögliches Monitoring}
% 
% \begin{itemize}
%   \item Vitale Bedrohung beurteilen. {\TxBeiVit} Behandlung
%   wie schweres Lungenöden (\prettyref{m:lungenoedem-schwer})
%   \item Lagerung: Oberkörper hoch
%   \item Beengende Kleidung öffnen
%   \item Bestmögliches Monitoring
% \end{itemize}
% }

\MassnahmenMK{Spezielle Maßnahmen}{Lungenödem,
schwer}{m:lungenoedem-schwer}
{
\EE{Taktik}: \Speech{\Ca{Vitale Bedrohung}
Belastung minimieren, rasche
medikamentöse Therapie (veranlassen); ärztliche Therapie
notwendig}


\Cave{Patienten mit Lungenödem und brodelndem Atemgeräusch oder ABCD-Problem 
sind grundsätzlich notarztpflichtig. 

Auch bei kurzer Transportzeit hat die Stabilisierung des Patienten vor Ort Vorrang.

Bereits der Transport in das Fahrzeug kann gefährlich sein!}
}{}
% \MassnahmenNG{Spezielle Maßnahmen}{Lungenödem,
% schwer}{\label{m:lungenoedem-schwer}}
% {
% \EE{Taktik}: \Speech{\Ca{Vitale Bedrohung}
% Belastung minimieren, rasche
% medikamentöse Therpie (veranlassen); ärztliche Therapie
% notwendig}
% 
% \begin{itemize}
%   \item \TxMassVit
%   \begin{itemize}
%     \item Lagerung: Oberkörper hoch, \E{wenn möglich Beine
%     herabhängen lassen}
% 
%     Striktes \EE{Bewegungsverbot!}
%     \item  Notarzt!
%     \item Reanimationsbereitschaft!
%   \end{itemize}
%   \item Beengende Kleidung öffnen
%   \item \Z{Evtl. Gabe von Nitroglycerin-Spray durch einen
%   Notfallsanitäter}
% \end{itemize}
% }




\section{Hyperventilationssyndrom und -tetanie}
\label{topic:hyperventilationstetanie}
\label{topic:hyperventailationssyndrom}
\index{Hyperventilation!-stetanie}
\index{Hyperventilation!-ssyndrom}
\index{Hyperventilation!psychisch bedingte --}



\Layout{0}{{\TxBeschreibung} und Ursachen}{%
\DefinitionInPara{Als \T{Hyperventilationssyndrom} wird die 
übermäßig tiefe und schnelle Atmung \E{ohne körperliche Ursache} 
bezeichnet.}
Durch die Atmung wird der Säure-Basen-Haushalt des Körpers
wesentlich beeinflusst
(\prettyref{topic:saeure-basen-haushalt}). Atmet ein sonst
gesunder Mensch übermäßig tief und schnell, so wird sein
Blut-pH sehr rasch basisch. Dadurch geraten die
\E{Elektrolyte} im Blut in ein Ungleichgewicht, es kommt zu
Symptomen.

Das Hyperventilationssyndrom ist meistens psychisch bedingt
und Folge einer \EE{psychischen Belastungsreaktion}
bzw. Folge einer Panikattacke. Auch im Rahmen von
(Pop-)Konzerten und ähnlichen Veranstaltungen sind oft
Patienten mit Hyperventilationstetanie anzutreffen.
\bigskip
}{
\begin{itemize}
  \item Tiefe, schnelle Atmung bei sonst gesundem Menschen
  \item Blut-pH wird basisch $\rightarrow$
  Elektrolytstörung
  \item Meist psychisch bedingt (Psych. Belastungsreaktion, Panikattacke, Prüfungen, Veranstaltungen)
%   \begin{itemize}
%     \item 
%     %   \item Lungenembolie
%     %   \item diabetische Stoffwechselentgleisung (Ketoazidose, kompensatorisch)
%   \end{itemize}
\end{itemize}
}{}{}{}

\Fazit{Beim Hyperventilationssysdrom ist die Atmung \EE{ohne 
körperliche Ursache} gesteigert!}

\Layout{0}{\TxSymptome}
{%
Die Patienten beklagen \EE{Atemnot}, obwohl die
\ce{O2}-Zufuhr nicht beeinträchtigt und das Blut maximal
gesättigt ist (Sp\ce{O2} von 100\PC*). Sie haben eine
\EE{tiefe und schnelle Atmung} (\E{Tachypnoe}).
Durch das Elektrolyt-Ungleichgewicht kommt es zuerst zu
\EE{Schwindelgefühlen}, etwas später zu einem \EE{Kribbeln} in den
Fingern. In weiterer Folge zeigen sich tonische \EE{Krämpfe} 
(\T{Hyperventilationstetanie}).
Klassisch ist die sog.
\T{Pfötchenstellung}: Die Arme sind
angezogen und Handhaltung erinnert an Pfoten eines Tieres.
Wird die Hyperventilation nicht rechtzeitig beendet, kann es
zur Eintrübung bis hin zur Bewusstlosigkeit kommen.
}{
\begin{itemize}
  \item Subjektiv: \emph{\protect\enquote{Atemnot}}, Sp\ce{O2} von 100\PC* 
  \item Tiefe, schnelle Atmung
  \item Schwindel, Bewusstseinsstörungen
  \item Kribbeln in den Fingern
  \item Tonische Krämpfe (\protect\enquote{Pfötchenstellung})
\end{itemize}
}{}{}{}



\Layout{0}{Differentialdiagnosen}
{%
Die psychisch bedingte Hyperventilation muß unbedingt von
der Hyperventilation aufgrund anderer Ursachen abgegrenzt
werden. Bei \EE{Ateminsuffizienz} oder \EE{durch
Erkrankungen verursachte Atemnot} kommt es natürlich auch
oft zu einer (kompensatorischen) Hyperventilation
(Herzinfarkt, Lungenembolie, Pneumothorax, Asthma-Anfall,
\ldots).

Bei einer stoffwechselbedingten \EE{Übersäuerung} (Azidose)
kommt es zu einer Hyperventilation um \ce{CO2} abzuatmen
(dadurch wird Säure im Körper abgebaut). Die
\E{Kussmaul'sche Atmung} beim \EE{diabetischen Koma} wäre
ein typisches Beispiel dafür
(\prettyref{topic:koma-diabetisches}).

}{
\begin{itemize}
  \item Atemnot mit körperlicher Ursache
  \item Ateminsuffizienz
  \item Stoffwechselbedingte Azidose (z.\,B. Diabetisches Koma)
\end{itemize}
}{}{}{}


% \paragraph{\TxSymptome}
% 
% \begin{itemize}
% \item Unruhe, Aufregung, evtl. Angstgefühle, Panik
% \item Tachypnoe mit Alkalose $\rightarrow$ durch Änderung
% des pH-Werts Änderung des Kalziumspiegels im Blut, dadurch
% \begin{itemize}
%     \item Kribbeln, 
% \item Krämpfe (Gesichtsmuskulatur),  Pfötchenstellung
% \end{itemize}
% \end{itemize}

\MassnahmenMK{Spezielle
Maßnahmen}{Hyperventilationssyndrom}{m:Hyperventilationssysndrom}
{ \EE{Taktik}:
\Speech{Beruhigen, Differentialdiagnosen ausschließen, symptomatische Therapie}
}{}
% \MassnahmenNG{Spezielle
% Maßnahmen}{Hyperventilationssyndrom}{\label{m:Hyperventilationssysndrom}}
% { \EE{Taktik}:
% \Speech{Beruhigen, Differentialdiagnosen ausschließen, symptomatische Therapie}
% 
% \begin{itemize}
%   \item Psychische Betreuung: Patient beruhigen, mit ruhiger, langsamer Stimme
%   mit dem Patienten sprechen
%   \item Kommandoatmung
%   \item \ce{CO2}-Rückatmung: In ein Plastiksackerl über kurze Zeit ein- und
%   ausatmen lassen. Ebenso eignet sich eine \EE{\ce{O2}-Maske mit Reservoir} oder ein (puderfreier) 
%   Handschuh.
%   \item Grundsätzlich sollte kein \ce{O2}  verabreicht werden. Aus
%   psychologischen Gründen kann es jedoch sinnvoll sein \EE{\ce{O2} in geringer
%   Dosierung (1\LPerMin*{})} über eine \ce{O2}-Maske zu
%   verabreichen.
%   \item Wenn nicht erfolgreich und Krämpfe bestehen bleiben bzw.
%   Bewusstseinsstörungen auftreten: \TxMassVit 
% \end{itemize}
% }




\ChapterEnd
